\documentclass[11pt,aspectratio=169,xcolor=dvipsnames]{beamer}

\usetheme[
  outer/progressbar=foot,
]{metropolis}

\usepackage{amsmath}
\usepackage{amssymb}
\usepackage{mathtools}
\usepackage{booktabs}
\usepackage{xcolor}
\usepackage{tikz}
\usepackage{fontspec}
\usepackage{polyglossia}
\setmainlanguage{polish}

\IfFontExistsTF{Fira Sans Light}{
  \setsansfont[BoldFont={Fira Sans}]{Fira Sans Light}
}{
  \setsansfont{Arial}
}
\IfFontExistsTF{Fira Mono}{
  \setmonofont{Fira Mono}
}{
  \setmonofont{Menlo}
}

\setbeamercolor{palette primary}{bg=BlueViolet,fg=white}
\setbeamercolor{palette secondary}{bg=BlueViolet,fg=white}
\setbeamercolor{title separator}{bg=Red,fg=BlueViolet}
\setbeamercolor{progress bar}{bg=OrangeRed,fg=BlueViolet}
\beamertemplatenavigationsymbolsempty

\newcommand{\rank}{\operatorname{rank}}

\title{Kointegracja i VECM}
\subtitle{ASC 2026}
\author{Piotr Żoch}
\date{}

\begin{document}

\begin{frame}[plain]
  \vspace{1.2cm}
  {\usebeamerfont{title}\usebeamercolor[fg]{title}Kointegracja i VECM\par}
  \vspace{0.35cm}
  {\usebeamerfont{subtitle}\usebeamercolor[fg]{subtitle}ASC 2026\par}
  \vspace{0.9cm}
  {\usebeamerfont{author}Piotr Żoch\par}
\end{frame}

\begin{frame}{Plan}
  \begin{itemize}
    \item Regresja pozorna
    \item Kointegracja
    \item Testowanie kointegracji
    \item ECM
    \item VECM
  \end{itemize}
\end{frame}

\section{Kointegracja}

\begin{frame}{Regresja pozorna}
  \small
  Modelowanie zależności między niepowiązanymi zmiennymi niestacjonarnymi może prowadzić do regresji pozornej:
  \begin{align*}
    x_t &= x_{t-1} + \varepsilon_t,\\
    y_t &= y_{t-1} + \nu_t.
  \end{align*}

  W regresji
  \[
    y_t = \alpha + \beta x_t + u_t
  \]
  często otrzymujemy wysokie $R^2$ i istotne statystyki $t$, mimo że zmienne są niezależne.

  \begin{itemize}
    \item Typowa reakcja: różnicowanie.
    \item Problem: różnicowanie może wyrzucić informację o długookresowej równowadze.
  \end{itemize}
\end{frame}

\begin{frame}{Wspólny trend stochastyczny}
  Zmienne mogą być niestacjonarne, ale mieć wspólny trend:
  \begin{align*}
    x_t &= z_t + \varepsilon_t,\\
    y_t &= z_t + \nu_t,\\
    z_t &= z_{t-1} + \zeta_t.
  \end{align*}

  Wtedy
  \[
    x_t - y_t = \varepsilon_t - \nu_t.
  \]

  \begin{itemize}
    \item $x_t$ i $y_t$ są $I(1)$.
    \item Różnica $x_t-y_t$ może być $I(0)$.
    \item Krótkookresowo zmienne mogą się od siebie oddalać, ale długookresowo poruszają się razem.
  \end{itemize}
\end{frame}

\begin{frame}{Kointegracja}
  Niech $x_t=(x_{1t},\ldots,x_{nt})'$ będzie wektorem zmiennych $I(d)$.

  Jeśli istnieje niezerowy wektor $\beta=(\beta_1,\ldots,\beta_n)'$, taki że
  \[
    \beta'x_t \sim I(d-b), \qquad b>0,
  \]
  to zmienne są skointegrowane stopnia $(d,b)$.

  \begin{itemize}
    \item W praktyce najczęściej: $d=b=1$.
    \item Wtedy kombinacja liniowa zmiennych $I(1)$ jest stacjonarna.
    \item $\beta$ nazywamy wektorem kointegrującym.
  \end{itemize}
\end{frame}

\begin{frame}{Wektor i rząd kointegracji}
  \begin{itemize}
    \item Kointegracja dotyczy kombinacji liniowej.
    \item Jeśli $\beta$ jest wektorem kointegrującym, to $\lambda\beta$ dla $\lambda\neq 0$ też nim jest.
    \item Dlatego zwykle normalizujemy jeden parametr, np. przyjmujemy współczynnik przy $y_t$ równy 1.
    \item Dla $n$ zmiennych może istnieć najwyżej $n-1$ liniowo niezależnych wektorów kointegrujących.
    \item Liczba takich wektorów to rząd kointegracji.
  \end{itemize}
\end{frame}

\begin{frame}{Długookresowa relacja}
  Jeśli $y_t$ i $x_t$ są skointegrowane, to w regresji
  \[
    y_t = \alpha + \beta x_t + u_t
  \]
  reszta $u_t$ jest stacjonarna.

  \begin{itemize}
    \item Interpretacja: $u_t$ to odchylenie od długookresowej równowagi.
    \item Gdy $y_t$ oddala się od $\alpha+\beta x_t$, system ma tendencję do powrotu.
    \item Estymator MNK relacji długookresowej jest super-zgodny, ale standardowe rozkłady statystyk nie działają bezpośrednio.
  \end{itemize}
\end{frame}

\begin{frame}{Kointegracja a różnicowanie}
  Jeśli zmienne są skointegrowane, regresja na różnicach:
  \[
    \Delta y_t = \gamma \Delta x_t + e_t
  \]
  opisuje tylko zależność krótkookresową.

  \begin{itemize}
    \item Tracimy informację o poziomie równowagi.
    \item Możemy otrzymać obciążony obraz zależności długookresowej.
    \item Potrzebujemy modelu, który zawiera:
      \begin{itemize}
        \item krótkookresową dynamikę w różnicach,
        \item korektę odchylenia od równowagi.
      \end{itemize}
  \end{itemize}
\end{frame}

\section{Testowanie kointegracji}

\begin{frame}{Testowanie kointegracji}
  Podstawowy schemat:
  \begin{enumerate}
    \item Sprawdź stopień integracji zmiennych.
    \item Jeśli zmienne są $I(0)$, kointegracja nie jest potrzebna.
    \item Jeśli zmienne są zintegrowane różnego stopnia, nie mogą być skointegrowane w standardowym sensie.
    \item Jeśli zmienne są $I(1)$, testuj stacjonarność relacji długookresowej.
  \end{enumerate}

  \vspace{0.3cm}
  Wybór testu:
  \begin{itemize}
    \item dwie zmienne: Engle--Granger,
    \item wiele zmiennych i możliwy rząd kointegracji większy niż 1: Johansen.
  \end{itemize}
\end{frame}

\begin{frame}{Test Engle--Granger}
  Krok 1: oszacuj relację długookresową:
  \[
    y_t = \alpha + \beta x_t + u_t.
  \]

  Krok 2: testuj pierwiastek jednostkowy w resztach:
  \[
    \Delta \hat{u}_t
      = \rho \hat{u}_{t-1}
      + \sum_{i=1}^{k}\phi_i \Delta \hat{u}_{t-i}
      + e_t.
  \]

  \[
    H_0:\rho=0
    \qquad
    H_1:\rho<0.
  \]

  Odrzucenie $H_0$ sugeruje kointegrację.
\end{frame}

\begin{frame}{Test Engle--Granger -- uwagi}
  \begin{itemize}
    \item Nie używamy zwykłych wartości krytycznych z testu ADF.
    \item Reszty $\hat{u}_t$ są skonstruowane po estymacji $\alpha$ i $\beta$.
    \item Wartości krytyczne zależą od:
      \begin{itemize}
        \item liczby zmiennych,
        \item deterministycznych składników w relacji długookresowej,
        \item trendu lub dryfu w zmiennych.
      \end{itemize}
    \item Test Engle--Granger zakłada jeden wektor kointegrujący i jest wrażliwy na wybór zmiennej po lewej stronie.
  \end{itemize}
\end{frame}

\begin{frame}{Test Johansena}
  Test Johansena wychodzi od reprezentacji VECM:
  \[
    \Delta y_t
      = \Pi y_{t-1}
      + \sum_{i=1}^{p-1}\Gamma_i \Delta y_{t-i}
      + \Phi D_t
      + \varepsilon_t.
  \]

  Kluczowy obiekt:
  \[
    r = \rank(\Pi).
  \]

  \begin{itemize}
    \item $r=0$: brak kointegracji.
    \item $0<r<n$: istnieje $r$ relacji kointegrujących.
    \item $r=n$: zmienne są stacjonarne w poziomach.
  \end{itemize}
\end{frame}

\begin{frame}{Test Johansena -- intuicja}
  Jeśli $0<r<n$, macierz $\Pi$ można rozłożyć jako
  \[
    \Pi = \alpha\beta',
  \]
  gdzie:
  \begin{itemize}
    \item $\beta$ zawiera wektory kointegrujące,
    \item $\beta'y_{t-1}$ to odchylenia od równowagi długookresowej,
    \item $\alpha$ opisuje tempo dostosowania poszczególnych zmiennych do tych odchyleń.
  \end{itemize}

  \vspace{0.3cm}
  Testy Johansena sprawdzają kolejne hipotezy o rzędzie $r$.
\end{frame}

\section{ECM i VECM}

\begin{frame}{Model korekty błędem}
  Twierdzenie Grangera o reprezentacji:

  jeśli $y_t$ i $x_t$ są $I(1)$ oraz skointegrowane, to zależność między nimi można przedstawić jako ECM:
  \[
    \Delta y_t
      = \kappa
      + \sum_{i=0}^{K-1}\gamma_i\Delta x_{t-i}
      + \sum_{i=1}^{L-1}\theta_i\Delta y_{t-i}
      + \delta (y_{t-1}-\alpha-\beta x_{t-1})
      + \varepsilon_t.
  \]

  \begin{itemize}
    \item Część w różnicach: dynamika krótkookresowa.
    \item Nawias: odchylenie od równowagi długookresowej.
  \end{itemize}
\end{frame}

\begin{frame}{Interpretacja ECM}
  Niech
  \[
    e_t = y_t - \alpha - \beta x_t.
  \]

  W ECM:
  \[
    \Delta y_t = \kappa + \delta e_{t-1} + \cdots + \varepsilon_t.
  \]

  \begin{itemize}
    \item $e_{t-1}$ mierzy odchylenie od równowagi.
    \item $\delta$ opisuje tempo powrotu do równowagi.
    \item Zwykle oczekujemy $\delta<0$ w równaniu dla $\Delta y_t$.
    \item Przybliżony okres połowicznego wygaśnięcia:
      \[
        -\frac{\ln 2}{\ln(1+\delta)}.
      \]
  \end{itemize}
\end{frame}

\begin{frame}{Od VAR do VECM}
  Startujemy od VAR($p$):
  \[
    y_t = A_1y_{t-1} + \ldots + A_py_{t-p} + \varepsilon_t.
  \]

  Po odjęciu $y_{t-1}$ i przekształceniu:
  \[
    \Delta y_t
      = \Pi y_{t-1}
      + \sum_{i=1}^{p-1}\Gamma_i\Delta y_{t-i}
      + \varepsilon_t,
  \]
  gdzie
  \[
    \Pi = A_1 + \ldots + A_p - I.
  \]

  \begin{itemize}
    \item To jest VECM: vector error correction model.
  \end{itemize}
\end{frame}

\begin{frame}{VECM}
  Gdy zmienne są skointegrowane:
  \[
    \Pi = \alpha\beta'.
  \]

  \[
    \Delta y_t
      = \alpha\beta'y_{t-1}
      + \sum_{i=1}^{p-1}\Gamma_i\Delta y_{t-i}
      + \Phi D_t
      + \varepsilon_t.
  \]

  \begin{itemize}
    \item $\beta'y_{t-1}$: relacje długookresowe.
    \item $\alpha$: prędkości dostosowania.
    \item $\Gamma_i$: krótkookresowa dynamika.
    \item $D_t$: stała, trend, sezonowość, zmienne interwencyjne.
  \end{itemize}
\end{frame}

\begin{frame}{VAR, różnice czy VECM?}
  \begin{center}
    \begin{tabular}{lll}
      \toprule
      Dane & Relacje długookresowe & Model \\
      \midrule
      $I(0)$ & nie dotyczy & VAR w poziomach \\
      $I(1)$ & brak kointegracji & VAR w różnicach \\
      $I(1)$ & kointegracja & VECM \\
      \bottomrule
    \end{tabular}
  \end{center}

  \vspace{0.4cm}
  \begin{itemize}
    \item VAR w poziomach może być dobry prognostycznie, ale wymaga ostrożnej inferencji.
    \item VECM łączy krótkookresową dynamikę z długookresową równowagą.
  \end{itemize}
\end{frame}

\begin{frame}[fragile]{R: podstawowy schemat}
\tiny
\begin{verbatim}
library(vars)
library(urca)

# 1. stopien integracji
summary(ur.df(y[, "m"], type = "drift", lags = 4))
summary(ur.df(y[, "p"], type = "drift", lags = 4))

# 2. test Johansena i estymacja VECM
jo <- ca.jo(y, type = "trace", ecdet = "const", K = 4)
summary(jo)
vecm <- cajorls(jo, r = 1)

# 3. VAR odpowiadajacy VECM, np. do IRF
var_from_vecm <- vec2var(jo, r = 1)
\end{verbatim}
\end{frame}

\begin{frame}{Najważniejsze wnioski}
  \begin{itemize}
    \item Regresja niestacjonarnych zmiennych może być pozorna.
    \item Kointegracja oznacza, że pewna kombinacja liniowa zmiennych $I(1)$ jest stacjonarna.
    \item Test Engle--Granger działa dla prostej relacji długookresowej.
    \item Test Johansena pozwala badać rząd kointegracji w układzie wielu zmiennych.
    \item VECM łączy krótkookresową dynamikę z korektą błędem.
  \end{itemize}
\end{frame}

\end{document}
